\chapter{From bytes to an image file}
\MARSChapterHeading{From bytes to an image file}
\label{revision-lab:files}

Use the complete \texttt{ppm-codec.asm} in \texttt{revision/labs/files}.
The main chapter establishes the restricted P6 format, caller-owned intervals
and file-service interfaces. This lab follows an image across those boundaries.
You should finish with an actual image file and an explanation of why its bytes
are correct. A viewer accepting a file is useful evidence, but it does not check
saved registers, untouched memory or an error path.

The native student archive is \texttt{029-binary-ppm-0.1.0-student.zip}.
Choose \texttt{Tools > Import MARS Lab...}, select the archive, and continue it
in Learn. Answer predictions in the inline choices beside the source pointer;
ordinary Run and Debug provide the execution observations. You can also work
through the written activities below without the guide. Keep a clean worked
source and use separate copies for deliberate faults. The complete implementation
provides helpers for the independent exercise; you need not rebuild every
function before producing a file.

The package includes complete variations under \texttt{worked/}. Their routine
bodies are unchanged; each variation selects a mode or a boundary input. The
pointer in Learn identifies the next requested operation. The selected Debug
frame identifies the instruction at which execution is actually paused. Use
one line breakpoint at a time and resume to normal exit after the last
observation before opening another variation. At the final exit syscall,
\texttt{\$v0} is ten, so inspect the caller's stored status to recover the
earlier function result.

\section{Read the default image without opening a file}

Keep \texttt{mode} zero. Locate \texttt{source}, \texttt{byte\_count},
\texttt{metadata}, \texttt{status} and \texttt{count}. There are 23 selected
source bytes, although the allocation is larger. Read the eleven-byte ASCII
header and account for the remaining twelve bytes. Predict width, height,
raster offset and pixel count before starting Debug.

Break at \texttt{parse\_return}. Expect status zero in \texttt{\$v0} and
four pixels in \texttt{\$v1}. The descriptor holds width two and height two
in its first two words, then offset eleven in its doubleword. The descriptor
contains an offset rather than the address of a pixel. The caller's
\texttt{status} and \texttt{count} objects still hold their initial $-1$ values:
this breakpoint precedes the caller's publication at \texttt{results\_ready}.
Do not confuse those objects with the live function results.

In a fresh copy, select only 22 bytes. Do not alter the backing allocation.
The header still describes four pixels, but the selected interval is missing
one sample. Predict the error and which metadata bytes may change. At return,
status is one, length zero, and the metadata retains its hexadecimal
\texttt{a5} sentinel bytes. Repeat with selected count 24. An additional mapped
byte also makes this single-image extent invalid. Explain why accepting a
plausible header would not be enough for either case.

\section{Put punctuation in the raster}

Make a one-by-one fixture with header \texttt{P6\textbackslash n1 1\textbackslash n255\textbackslash n}
and sample bytes ten, 35 and zero. Its assembly initializer is:
\begin{verbatim}
source: .byte 80,54,10,49,32,49,10,50,53,53,10,10,35,0
    .byte -91:12562
\end{verbatim}
Set \texttt{byte\_count} to fourteen. The filler keeps the source arena at
12576 bytes so that later objects stay at the same addresses; filler does not
belong to the selected file. Before running, identify which LF ends the header
and which LF is the first red sample. Identify the byte that resembles a
comment marker but is actually green.

At \texttt{parse\_return}, the descriptor must contain one, one and eleven.
The parser must not skip either raster sample. The native observation below
shows that descriptor for this exact fixture.

\begin{figure}[htbp]
\centering
\begin{tikzpicture}[alt={The actual light-theme Memory view after parsing the one-pixel punctuation fixture. Its sixteen descriptor bytes encode width one, height one and raster offset eleven. Raster samples LF, number-sign and zero have not been consumed as header syntax.}]
\node[anchor=north west,inner sep=0,font=\small\bfseries] (title) at (0,0) {Width, height and raster offset};
\node[anchor=north west,inner sep=0] (a) at ([yshift=-2mm]title.south west) {\MARSRevisionCrop{revision/captures/ppm/ppm-parser-boundary.png}{0.99\linewidth}{0.03125}{0.31555556}{0.58203125}{0.60888889}};
\begin{scope}[shift={(a.south west)},x={($(a.south east)-(a.south west)$)},y={($(a.north west)-(a.south west)$)}]
\draw[black,line width=.8pt] (.20,.3088) rectangle (.965,.6176);\draw[black,line width=.8pt] (.58,.3088)--(.58,.6176);
\end{scope}
\end{tikzpicture}
\caption{The first eight bytes hold width and height; the next eight hold offset eleven. The parser has accepted all fourteen selected file bytes. Native LaTeX supplies the title, crop and outlines.}
\label{fig:revision-ppm-parser}
\end{figure}


Switch to mode one in another copy of the same fixture. Break at
\texttt{decode\_return}. Expect status zero and packed dimensions
\texttt{0x0000000100000001}. The first framebuffer word is
\texttt{0x000a2300}; in little-endian byte view it appears as
\texttt{00 23 0a 00}. Explain every byte before interpreting the color.

For a diagnosis exercise, imagine a parser that invokes its general gap skipper
after Maxval. Follow it across the two consecutive LF bytes and the following
\texttt{\#}. Which selected bytes would it misclassify? Describe the missing
pixel boundary rather than merely saying that the resulting image looks wrong.

\section{Serialize dimensions that have different digit counts}

Return to an unchanged worked source, choose mode two and set
\texttt{dimensions} to decimal 42949672969. This packs height ten in the high
word and width nine in the low word. The framebuffer allocation has sufficient
capacity; its initialized first four words are red, green, blue and white, and
later words contain the fixture sentinel pattern. This is a header experiment,
not a claim that the remaining pixels form a designed picture.

Derive the canonical header \texttt{P6\textbackslash n9 10\textbackslash n255\textbackslash n}.
Its length is twelve: one dimension needs a second decimal digit. The raster
needs $9\cdot10\cdot3=270$ bytes, so success returns 282 bytes in
\texttt{\$v1}. Break at \texttt{encode\_return} and inspect the destination.
The first red sample follows the twelfth header byte immediately.

\begin{figure}[htbp]
\centering
\begin{tikzpicture}[alt={The actual light-theme Memory view after encoding width nine and height ten. Its first twelve bytes encode P6 newline 9 space 10 newline 255 newline. The next bytes begin the red pixel with ff 00 00.}]
\node[anchor=north west,inner sep=0,font=\small\bfseries] (title) at (0,0) {Twelve header bytes, then RGB samples};
\node[anchor=north west,inner sep=0] (a) at ([yshift=-2mm]title.south west) {\MARSRevisionCrop{revision/captures/ppm/ppm-encoded-header.png}{0.99\linewidth}{0.03125}{0.31555556}{0.58203125}{0.60888889}};
\begin{scope}[shift={(a.south west)},x={($(a.south east)-(a.south west)$)},y={($(a.north west)-(a.south west)$)}]
\draw[black,line width=.8pt] (.20,.3088) rectangle (.77,.6176);\draw[black,line width=.8pt,dashed] (.77,.3088) rectangle (.965,.6176);
\end{scope}
\end{tikzpicture}
\caption{The outlined twelve-byte header ends before the red pixel. Different decimal digit counts move the raster start without changing its three-byte pixel stride. Native LaTeX supplies the title, crop and outlines.}
\label{fig:revision-ppm-header}
\end{figure}


Change capacity to 281 in a separate copy. Predict status two, zero length and
an unchanged entire destination, including the header positions. The encoder
calculates and checks the complete extent before emitting either text or samples.
Then restore exact capacity 282 and verify success without a trailing NUL.
Why would checking only the three-byte-per-pixel raster requirement be too weak?

\section{Run the complete pipeline on actual files}

The worked mode-five caller loads, decodes, encodes and saves. The following
command runs a source copy in a newly created directory and verifies the output
bytes. Run it from the public book repository, replacing the two tool paths.
Choose an output directory that does not exist; the helper refuses to reuse one.
\begin{verbatim}
python3 tools/run_revision_ppm.py \
  --java /path/to/java --cli /path/to/cli.jar \
  --output /tmp/ppm-first-run
\end{verbatim}
The helper writes \texttt{input.ppm}, prepares \texttt{pipeline.asm}, executes
MARS and checks \texttt{output.ppm}. The default image is two by two: red, green,
blue and white. The output has 23 bytes. It replaces a deliberately longer old
output, testing that a successful save truncates it rather than leaving stale
trailing bytes. Console silence is expected; the result is a file.

Keep the resulting directory. Its receipt records the executed command and
hashes, and its output file can be opened with a PPM-capable image viewer. At
this tiny size, use magnification without smoothing when identifying individual
pixels. A successful receipt checks the exact canonical bytes; it does not
measure performance or certify every PPM variant.

Repeat in another new directory with \texttt{--width 9 --height 10 --commented}.
The input contains a comment adjacent to the magic and uses a tab between the
dimensions, both accepted by our local parser. The canonical output replaces
that header and retains the samples. Some external readers, including the
FFmpeg version used for this draft's verification, do not recognize that
comment-adjacent-magic input. The generated canonical output was independently
decoded successfully. Keep local acceptance rules and a particular viewer's
behavior distinct.

\section{Account for partial transfers and ownership}

A file descriptor belongs to the function that obtained it until it is closed.
It is neither a path pointer nor necessarily the integer three. A read request
is a maximum; a positive result may be smaller. After accepting $k$ bytes, the
loader asks for the remaining capacity at \texttt{buffer+k}. Zero means EOF;
a negative result means failure. Explain why a positive short read alone does
not establish EOF.

\begin{figure}[htbp]
\centering
\begin{tikzpicture}[alt={The actual light-theme Memory view at the complete file pipeline return. The first sixteen framebuffer bytes hold red, green, blue and white as little-endian 00RRGGBB words. Separate operation evidence records forced short transfers and closed descriptor nineteen.}]
\node[anchor=north west,inner sep=0,font=\small\bfseries] (title) at (0,0) {Four decoded words after load, encode and save};
\node[anchor=north west,inner sep=0] (a) at ([yshift=-2mm]title.south west) {\MARSRevisionCrop{revision/captures/ppm/ppm-pipeline-pixels.png}{0.99\linewidth}{0.03125}{0.31555556}{0.58203125}{0.60888889}};
\begin{scope}[shift={(a.south west)},x={($(a.south east)-(a.south west)$)},y={($(a.north west)-(a.south west)$)}]
\draw[black,line width=.8pt] (.20,.3088) rectangle (.965,.6176);
\end{scope}
\end{tikzpicture}
\caption{The framebuffer retains red, green, blue and white after the complete pipeline. This Debug run uses the documented virtual file fixture; actual host-file verification is separate. Native LaTeX supplies the title, crop and outlines.}
\label{fig:revision-ppm-pixels}
\end{figure}


This Debug observation uses a deterministic virtual file installed by the
separate authoring plugin. It begins with the same 23 input bytes and forces
reads and writes of at most seven bytes. Descriptor 19 is opened and closed for
loading, then opened and closed for saving. The output bytes match the input's
canonical image. The screenshot establishes the decoded words at that checkpoint;
the archived operation receipt establishes the forced transfers and closes.
The earlier command separately exercises actual host files.

When an input fills the offered capacity exactly, the loader still needs to
know whether the file has another byte. Its one-byte probe uses private stack
storage, outside the caller's result interval. EOF accepts the full buffer;
one extra byte reports a capacity error. Drawing that extra byte into
\texttt{buffer+capacity} would write outside the offered interval even if more
physical memory happened to be mapped.

The saver validates the complete selected image before its truncating open.
Once open, it repeats writes until the selected count is accepted. A negative
return or a zero return for a nonempty request is an I/O failure; repeating the
same zero-progress request indefinitely cannot complete the file. The grader
has an explicit zero-progress control to test that case. A failure after an
open may leave a partial output file, but the descriptor must be closed and
the returned length must be zero. A malformed source, by contrast, must leave
an existing output file unopened and unchanged.

In your notes, draw separate timelines for an invalid source, a failure after
eleven accepted bytes, and a successful sequence of short writes. Mark the
first operation permitted to alter the output file. State what is guaranteed
about file contents and descriptor ownership on each return path.

\section{Connect the guided observations}

The first observations use \texttt{ppm-codec.asm} before its descriptor stores,
at the parser return and after caller publication. Record one line for each
boundary: the values derived so far, the memory objects that have changed,
and the value the caller may now use. A parser can have width two and height
two in registers while its caller-visible descriptor still consists of
\texttt{a5} bytes. Later, the descriptor can be valid while the caller's
\texttt{status} and \texttt{count} still contain their initial values. These
are different stages of one successful operation.

Follow the same argument through \texttt{truncated-raster.asm} and
\texttt{extra-raster-byte.asm}. Both reject the selected extent without
publishing metadata. Then use \texttt{raster-punctuation.asm} and
\texttt{decode-punctuation.asm} to locate the transition from header syntax
to samples. Derive the decoded word from the three \texttt{lbu} results before
reading its four stored bytes. A correct answer about the word's hexadecimal
value should also explain why its byte order differs from the file's channel
order.

The serialization observations use \texttt{encode-nine-by-ten.asm},
\texttt{capacity-short.asm} and \texttt{encode-maximum.asm}. Work out the
complete extent before checking the offered capacity. The first fixture
needs 282 bytes; the one-byte-short fixture may emit none of them. The
maximum fixture needs 12,301 bytes. In \texttt{count-bound.asm}, the parser
rejects a selected byte count above its documented limit before examining the
header. Keep a bound failure distinct from a format failure: they have
different causes and different status values.

The packaged \texttt{pipeline.asm} uses the supplied \texttt{worked/\allowbreak input.ppm}
and creates or replaces \texttt{worked/\allowbreak output.ppm}. Relative file paths are
resolved from this source directory. Its ordinary Run and Debug sessions use
host files; they do not install the separate authoring plugin's short-transfer
fixture. The returned descriptor is three in the supplied ordinary observation,
but the function must save and use whatever its open returns. The written
figure's descriptor 19 demonstrates why a fixed descriptor would be a bug.
The figure's forced seven-byte transfers remain a distinct deterministic test.

Keep the successful output before investigating the failures. In
\texttt{missing-input.asm}, the path is \texttt{absent-input.ppm}, which the
package does not supply. Check that it does not exist before expecting status
three. In \texttt{save-invalid.asm}, the selected source is 22 bytes long.
Expect format status one and count zero before the saver opens the output.
Compare the existing output bytes before and after; they must remain unchanged.
Creating or modifying a host file changes the experiment's input environment,
even when the assembly source is unchanged. Repeat the relevant observations
after such a change.

Learn checks the requested register states at Debug boundaries. Its ordinary
Run check establishes normal silent termination. Neither observation alone
proves the output file's bytes: make the explicit comparison and retain the
artifact. A file viewer adds another observation, while the standalone helper
below checks the canonical file against its independent host model. This is
why the lab asks for bytes, a visible image and an explanation rather than
merely an empty console or a completed progress marker.

Four independent assessment actions group the finite function corpus: 80
one-digit-width trials, 80 two-digit-width trials, six invalid dimensions and
the default caller. A fifth checks the complete caller's normal silent exit.
The many bounded data files are implementation details behind those five
actions. Equivalent correct solutions are accepted; the Debug observations of
the worked source do not prescribe the independent routine's temporary registers.
Use the graduated hints if your instructor has enabled them. Return to a failed
case with a specific prediction about the next read, write or returned value.

An edit invalidates execution passes for that source while retaining the
prediction history. Continue Downloaded MARS Lab resumes saved progress; it
should not be used to treat a result from an earlier source revision as current.
The independent checks use the active editor text, including unsaved edits.
The standalone script remains available for free writing, and the instructor
can control hint release, prediction prerequisites and assessment attempts.

\section{Write a complementary-color serializer}

Complete \texttt{ppm\_encode\_negative} in a copy of
\texttt{revision/labs/files/ppm-negative.asm}. Keep its supplied helpers and
caller. Arguments and failures match \texttt{ppm\_encode}: full framebuffer
pointer, packed dimensions, byte destination and unsigned capacity. Serialize
$255-x$ for each red, green and blue component; ignore the input word's high
byte. Preserve the source, saved registers, frame pointer and stack pointer.
Return status zero and the exact canonical byte count on success, otherwise
the specified error and zero count without output writes. Objects are disjoint.

For the default four colors, predict cyan, magenta, yellow and black. Write the
expected twelve raster bytes before running. The header remains the ordinary
PPM header: complementing it would destroy the file format. You may serialize
the complemented channels directly, or reuse the supplied encoder and transform
only the accepted raster after success. Neither design may temporarily change
the source or write rejected output and restore it later.

Use the public behavior checks:
\begin{verbatim}
java -jar /path/to/cli.jar grade \
  revision/labs/checks/ppm-negative.kts \
  /path/to/your/ppm-negative.asm
\end{verbatim}
The unchanged starter deliberately fails. The checks accept equivalent correct
implementations and test selected dimension, offset and capacity boundaries,
source preservation and caller state. Passing them does not prove every possible
pixel combination. Keep a high input byte other than zero, both one- and
two-digit dimensions, exact capacity and one byte less in your own explanation.

Finally run your completed routine through the actual file pipeline:
\begin{verbatim}
python3 tools/run_revision_ppm.py \
  --java /path/to/java --cli /path/to/cli.jar \
  --source /path/to/your/ppm-negative.asm --negative \
  --output /tmp/ppm-negative-run
\end{verbatim}
The helper selects your routine in a separate pipeline copy and checks the
canonical header and complemented raster. Keep the image and receipt, and
explain one prediction you corrected.
